\section{Introdução e Objetivos}

\subsection{Introdução}

Diversos fenômenos da natureza se repetem ao longo do tempo, e esses movimentos repetitivos recebem o nome de oscilações. Todo corpo que executa um movimento de vaivém em torno de uma posição de equilíbrio é denominado oscilador e, quando a força restauradora é proporcional ao deslocamento, tem-se um oscilador harmônico (NUSSENZVEIG, 2014).

No sistema massa-mola vertical, a posição de equilíbrio é aquela em que a força elástica da mola compensa o peso do corpo suspenso. Quando o corpo é afastado desse ponto e o meio é pouco dissipativo, como o ar, o sistema se comporta aproximadamente como um oscilador livre, com frequência angular natural $\omega_0=\sqrt{k/m}$, em que $k$ é a constante elástica da mola e $m$ é a massa suspensa (HALLIDAY; RESNICK; WALKER, 2016). Em um meio viscoso, como a água, a força de atrito viscoso faz com que a amplitude decaia exponencialmente no tempo e a frequência de oscilação diminua para um valor $\omega_1<\omega_0$. Por fim, quando uma força externa periódica de frequência angular $\Omega$ atua sobre o sistema, a massa passa a oscilar com essa frequência, e a amplitude torna-se função de $\Omega$, atingindo um máximo na chamada frequência de ressonância $\Omega_r$ (UNIVERSIDADE DE SÃO PAULO, 2026).

\subsection{Objetivos}

Este relatório tem como objetivo estudar o comportamento de um oscilador massa-mola vertical quanto à amplitude e à frequência das oscilações, em função da viscosidade do meio (ar e água), nas condições de oscilação livre, amortecida e forçada. Especificamente, busca-se:
\begin{itemize}
    \item determinar as frequências angulares de oscilação no ar ($\omega_0$) e na água ($\omega_1$) e a constante elástica $k$ da mola;
    \item determinar o fator de amortecimento $\gamma$ na água a partir do decaimento das amplitudes e verificar sua consistência com o modelo teórico;
    \item estudar o fenômeno da ressonância na oscilação forçada, determinando $\Omega_r$ no ar e na água e discutindo o efeito do amortecimento sobre a forma da curva de amplitude e sobre a posição do máximo;
    \item discutir as fontes de erro envolvidas em cada procedimento.
\end{itemize}
